% TODO proof-reading
\subsection{条件ジャンプ}

それでは条件ジャンプに進みましょう。

\begin{lstlisting}[style=customjava]
public class abs
{
	public static int abs(int a)
	{
		if (a<0)
			return -a;
		return a;
	}
}
\end{lstlisting}

\begin{lstlisting}
  public static int abs(int);
    flags: ACC_PUBLIC, ACC_STATIC
    Code:
      stack=1, locals=1, args_size=1
         0: iload_0       
         1: ifge          7
         4: iload_0       
         5: ineg          
         6: ireturn       
         7: iload_0       
         8: ireturn       
\end{lstlisting}

\TT{ifge}は\ac{TOS}の値が0以上の場合、オフセット7にジャンプします。

忘れてはいけないのは、任意の\TT{ifXX}命令は（比較される）値をスタックからポップすることです。

\TT{ineg}は\ac{TOS}の値を単純に否定します。

別の例：

\begin{lstlisting}[style=customjava]
	public static int min (int a, int b)
	{
		if (a>b)
			return b;
		return a;
	}
\end{lstlisting}